% Source: https://artofproblemsolving.com/community/c3h1908159_number_theory_problem
% Stable identity: aops-topic-1908159-q01
Let $r,b,g$ be the numbers of red, blue and green chameleons. In every meeting two of these numbers decrease by $1$ and the third increases by $2$. So each difference $r-b$, $b-g$, $g-r$ either stays the same or changes by $\pm3$ (for example, if $r$ and $b$ decrease, then $r-b$ is unchanged and $b-g$ changes by $-3$). Hence the remainders of $r-b$, $b-g$ and $g-r$ modulo $3$ never change.

At the start $r-b=-2$, $b-g=-2$ and $g-r=4$, and none of these is divisible by $3$. If at some moment all chameleons had the same colour, two of the three counts would be $0$, and their difference $0$ is divisible by $3$. This contradicts the invariant.

So the answer is $\boxed{\text{No}}$: the chameleons can never all have the same colour.
